\begin{abstract}
This report examines how floating-point precision, build semantics, solver
choice and hardware affect a tested ideal magnetohydrodynamics (MHD)
implementation using high-resolution shock-capturing (HRSC) methods. The
validation hierarchy separates reference error from cross-variant discrepancy
and stochastic arithmetic spread. It combines unit and property tests,
Brio--Wu numerical-reference comparison, two-dimensional three-grid checks, and
matched central processing unit (CPU) and graphics processing unit (GPU)
comparisons. Under the per-line positivity fallback, the Harten--Lax--van Leer
(HLL) and HLL with Discontinuities (HLLD) solvers kept density and pressure
finite and positive in every reported run. Across Orszag--Tang and
Kelvin--Helmholtz refinements, all 12
mean-$L_1$ density discrepancies between 32-bit (fp32) and 64-bit (fp64)
floating-point calculations remained below 3\% of the matched fp64 finest
adjacent-grid difference. This diagnostic relates discrepancy to discretisation
change; it establishes neither physical accuracy nor fp32 adequacy. In the
domain mean the difference amounted to between two and seven binary32
roundings on the coarse grids, but grew to several hundred by $512^2$, so
refinement widened the precision gap rather than closing it. Saved HLL CPU
and GPU states were bit-identical in the covered tests, although only because
device multiply--add contraction was disabled to match the host default:
restoring the compiler's own default removed the agreement in every pair and,
on Brio--Wu in fp32, moved the stored density further than halving the working
precision did. That agreement is therefore a property of the declared build
recipe, not of the hardware. The GPU was slower for small Brio--Wu and faster for larger
Orszag--Tang. Compiler floating-point semantics and HLLD branching produced
nonzero discrepancies in other matched variants. For Brio--Wu, the Monte Carlo
arithmetic (MCA) spread at virtual 24-bit precision (p24) was of the same order
as the deterministic fp32--fp64 difference. Robustness over the tested horizons
is therefore distinct from arithmetic insensitivity; reproducibility requires
complete records of the numerical path and execution environment.
\end{abstract}
